\section{Descripción del ámbito productivo en el que se realiza la PPII }

La Práctica Profesional II fue realizada en el Departamento de Neurociencia de la Facultad de Medicina de la Universidad de Chile. Específicamente, se llevó a cabo en la línea de investigación en neuromodulación y control motor, dirigida por el Dr. Rómulo Fuentes.\\

El departamento se dedica principalmente a la investigación científica, contando con diversas líneas en las que trabajan investigadores enfocados en distintas temáticas dentro del ámbito de la neurociencia.\\

El trabaj de práctica se desarrolló bajo la supervisión de la Dra. Andrea Colins y tuvo como objetivo el análisis de datos de actividad neuronal del ganglio pedal de ejemplares jóvenes y de mayor edad de Aplysia. Para ello, se aplican métodos estadísticos con el fin de comparar la actividad neuronal entre ambos grupos etarios.\\

Entre los insumos utilizados se incluye un computador con el rendimiento necesario para procesar grandes volúmenes de datos, una licencia estudiantil de MATLAB proporcionada por la Universidad de Chile y los registros neuronales, los cuales fueron facilitados por la Dra. Andrea Colins.\\

Los datos entregados corresponden a seis grabaciones en total: tres de Aplysia jóvenes y tres de Aplysia de mayor edad. Cada grabación contenía un archivo con la posición espacial de cada neurona, y otro con los tiempos en que cada una genera un spike.\\
